\begin{figure}[htp!]
  \centering
  \includegraphics[scale=0.22]{Esercizio13.2/Esercizio.jpg} 
\end{figure}
\textbf{A)} La risonanza $\Delta^+$ ha: $I=\frac{3}{2}$ e $I_3=+\frac{1}{2}$. Attraverso i coefficienti di Clebsch-Gordan è possibile andare a vedere in che stati può essere. Sicuramente sto facendo la composizione di $1\times1/2$ in quanto abbiamo un pione e un nucleone nello stato finale. Infatti ricordo che $N$ (nucleone) ha: $I=\frac{1}{2}$ e $I_3=\pm\frac{1}{2}$. Mentre il pione ha: $I=1$ e $I_3=\pm1,0$. Vedo che lo stato lo trovo in una combinazione di:
\begin{equation}
    \sqrt{\frac{1}{3}}|\pi^+n\rangle+\sqrt{\frac{2}{3}}|\pi^0p\rangle
\end{equation}
Il $BR$ è proporzionale all'elemento di matrice al quadrato quindi:
\begin{equation}
    \frac{\mathrm{BR}(\Delta^+\,\to\,\pi^+n)}{\mathrm{BR}(\Delta^+\,\to\,\pi^0p)}=\frac{1}{2}
\end{equation}
Quindi il decadimento in $\pi^+n$ avverà la metà delle volte rispetto al decadimento $\pi^0p$

\textbf{B)} Procedendo allo stesso modo si trova che $\Delta^0$ ha: $I=\frac{3}{2}$ e $I_3=-\frac{1}{2}$. Facendo riferimento ai coefficienti di Clebsch-Gordan:
\begin{equation}
    \sqrt{\frac{2}{3}}|\pi^0n\rangle+\sqrt{\frac{1}{3}}|\pi^-p\rangle
\end{equation}
E quindi:
\begin{equation}
    \frac{\mathrm{BR}(\Delta^0\,\to\,\pi^0n)}{\mathrm{BR}(\Delta^0\,\to\,\pi^-p)}=2
\end{equation}

\textbf{C)}  Gli stati di risonanza $N$ sono come i nucleoni ma in stato risonante quindi hanno: $I=\frac{1}{2}$ e $I_3=\pm\frac{1}{2}$. Per lo stato $N^+$ trovo che:
\begin{equation}
    \sqrt{\frac{2}{3}}|\pi^+n\rangle-\sqrt{\frac{1}{3}}|\pi^0p\rangle
\end{equation}
Quindi:
\begin{equation}
     \frac{\mathrm{BR}(N^+\,\to\,\pi^+n)}{\mathrm{BR}(N^+\,\to\,\pi^0p)}=2
\end{equation}

\textbf{D)} Per lo stato $N^0$:
\begin{equation}
    \sqrt{\frac{1}{3}}|\pi^0n\rangle-\sqrt{\frac{2}{3}}|\pi^-p\rangle
\end{equation}
Quindi:
\begin{equation}
    \frac{\mathrm{BR}(N^0\,\to\,\pi^0n)}{\mathrm{BR}(N^0\,\to\,\pi^-p)}=\frac{1}{2}
\end{equation}